\chapter*{Introducción}
\section*{Contextualización y antecedentes}

En el entorno empresarial contemporáneo, la transformación digital se ha convertido en un pilar fundamental para la supervivencia y competitividad de las organizaciones. No obstante, este proceso enfrenta barreras significativas, especialmente en las pequeñas y medianas empresas (PYMES), donde la escasez de recursos económicos y el desconocimiento de factores críticos limitan la adopción efectiva de tecnologías. Las PYMES a menudo operan con limitaciones financieras y brechas tecnológicas que restringen su capacidad de innovación.
En este contexto se sitúa la fábrica Confort Muebles, una empresa fundada en 2009 en Cuenca, Ecuador, dedicada a la fabricación y comercialización de muebles para el hogar y oficina, incluyendo diseños a medida. A pesar de su trayectoria y compromiso con la calidad, la empresa ha mantenido una gestión operativa basada tradicionalmente en procesos manuales para la administración de su información comercial y administrativa.

\section*{Planteamiento del problema}
La fábrica Confort Muebles enfrenta actualmente serias deficiencias en su gestión de proformas y procesos administrativos debido al uso predominante de registros manuales y archivos físicos independientes
. Esta situación genera una serie de inconvenientes críticos: errores recurrentes en los datos de precios y materiales, duplicidad de documentos, pérdida ocasional de información y una notable dificultad para realizar el seguimiento de clientes
.
El diagnóstico inicial revela que la elaboración de una sola proforma puede tardar entre 20 y 40 minutos, ya que los cálculos de costos de materiales, mano de obra y utilidad se realizan manualmente
. La falta de una base de datos centralizada impide el acceso rápido a la información, retrasa la atención al cliente y afecta la imagen institucional de la fábrica

\section*{Pregunta de investigación}
De qué manera la implementación de una aplicación web automatizada optimizará el control de proformas en la fábrica Confort Muebles?

\section*{Justificación}

\section{Objetivos}
Objetivo general
Implementar una aplicación web para la automatización del control de proformas en la fábrica Confort Muebles de la ciudad de Cuenca

Objetivos específicos:

Describir el flujo de información en los procesos de producción de la empresa Confort Muebles para identificar los requerimientos funcionales.

Diseñar la arquitectura tecnológica para la gestión de proformas, asegurando una estructura escalable y segura.

Implementar la aplicación web para la gestión de proformas de la empresa Confort Muebles.

Evaluar la integridad de la información de los procesos de la empresa Confort Muebles.
